\begin{table}[htbp]
\centering
\small
\caption{Contrastes de normalidad de la mortalidad}\label{tab:normalidad}
\begin{tabularx}{\linewidth}{X r r}
\toprule
Contraste & Estadístico & $p$ \\
\midrule
Shapiro-Wilk & \num{0.9903} & $<\num{0.001}$ \\
Kolmogorov-Smirnov-Lilliefors & \num{0.0280} & \num{0.001} \\
D'Agostino-Pearson K² & \num{127.3633} & $<\num{0.001}$ \\
Jarque-Bera & \num{269.8196} & $<\num{0.001}$ \\
Anderson-Darling & \num{4.0742} & \num{0.010} \\
$\chi^2$ de Pearson (clases equiprobables) ($k = 50$, 47 g.\,l.) & \num{94.2701} & $<\num{0.001}$ \\
Asimetría (contraste de D'Agostino) & \num{0.275} & $<\num{0.001}$ \\
Curtosis de exceso (contraste de Anscombe-Glynn) & \num{1.350} & $<\num{0.001}$ \\
\bottomrule
\end{tabularx}
\par\smallskip{\footnotesize\raggedright Anderson-Darling informa un $p$ interpolado en tablas (acotado entre 0,01 y 0,15). Con $n$ grande, cualquier desviación mínima resulta significativa: los contrastes se leen junto al gráfico Q-Q.\par}
\end{table}
